
\section*{关于本书 / About This Book}

\textbf{《工程本体论》} 系统讲述如何把本体论理论用于工程领域，覆盖概念化、
描述逻辑、RDF/RDFS/OWL、SPARQL、SHACL、推理、知识图谱以及
ontology-guided LLM/Agent。书中的概念、方法、案例叙事与习题是稳定的“石头”；
与这些内容对应的本体、能力问题（CQ）、形状、查询、案例、工程规则、执行合同、
版本与溯源收据，统一由 \textbf{Semantica} 承载和执行。

本目录是书稿，不是第二套语义运行库。删除或改写本目录中的文字不会绕过
Semantica 的包注册、能力声明、精确 oracle 与发布门禁；反过来，某个包场景运行
成功也只是在其声明边界内复算书中主张，不会自动把整章标记为完整或可发布。

\section*{唯一执行语义 / Sole Executable Semantics}

\begin{itemize}
\item 九章分别绑定 \texttt{vol1.ch01} 至 \texttt{vol1.ch09}；完整包名使用下文所列统一前缀。
\item RDF/OWL、CQ、SPARQL、SHACL、事实、规则、案例、合同和场景只在
Semantica 的 allowlisted built-in package 中保留可执行正本。
\item 本书中的代码框是教学节选；其历史来源路径用于溯源，构建时通过迁移账解析到
Semantica 包资产，不是可从本目录直接运行的副本。
\item Python、CLI 与 MCP 都必须委托同一个 \texttt{SemanticPackageRunner}；未知包、未知场景、
缺失资产、哈希不符或超出能力范围时 fail closed。
\item \texttt{run} 的场景 oracle 与 \texttt{verify} 的发布判定必须分开报告。当前各章 manifest 的
\texttt{release\_status} 均为 \texttt{blocked}，不得把教学片段或部分场景通过写成整章“全绿”。
\end{itemize}

\section*{书稿结构 / Book Structure}

\begin{codebox}
\codeline{ontology-engineering-book/}
\codeline{├──\ ch01-introduction/\ \ \ \ \ \ \ \ \ \ \ \ \ \#\ 哲学到工程的转向；包\ vol1.ch01}
\codeline{├──\ ch02-ontology-foundations/\ \ \ \ \ \#\ 基础理论；包\ vol1.ch02}
\codeline{├──\ ch03-ontology-methodology/\ \ \ \ \ \#\ 构建方法论；包\ vol1.ch03}
\codeline{├──\ ch04-ontology-languages/\ \ \ \ \ \ \ \#\ RDF/RDFS/OWL/SPARQL；包\ vol1.ch04}
\codeline{├──\ ch05-reasoning/\ \ \ \ \ \ \ \ \ \ \ \ \ \ \ \ \#\ 推理机制；包\ vol1.ch05}
\codeline{├──\ ch06-applications/\ \ \ \ \ \ \ \ \ \ \ \ \ \#\ 工程案例；包\ vol1.ch06}
\codeline{├──\ ch07-knowledge-graph/\ \ \ \ \ \ \ \ \ \ \#\ 知识图谱与\ SHACL；包\ vol1.ch07}
\codeline{├──\ ch08-ontology-llm/\ \ \ \ \ \ \ \ \ \ \ \ \ \#\ 本体\ \(\times\)\ LLM/Agent；包\ vol1.ch08}
\codeline{├──\ ch09-capstone-manufacturing/\ \ \ \#\ 制造业综合案例；包\ vol1.ch09}
\codeline{├──\ resources/\ \ \ \ \ \ \ \ \ \ \ \ \ \ \ \ \ \ \ \ \ \#\ 阅读与选型线索，不是自动信任清单}
\codeline{└──\ handbook/\ \ \ \ \ \ \ \ \ \ \ \ \ \ \ \ \ \ \ \ \ \ \#\ XeLaTeX\ 作者源、受控生成片段与成书\ PDF}
\end{codebox}
\color{black}

第 2、4、5、6、9 章的教学材料来自原书稿；第 1、3、7、8 章包含按全书脉络补写的
材料。所有已迁移材料在 Semantica manifest 中保留 \texttt{source\_anchor}、哈希与
\texttt{derivation}，以区分原文、结构化提取、适配器和新增执行合同。

\section*{九章执行落点 / Package Map}

下表以 \texttt{vol1.chNN} 缩写包名；其完整前缀均为
\texttt{semantica.chapter\_packages.}。

\begin{center}
\begin{tabularx}{\linewidth}{@{}lXlX@{}}
\toprule
\textbf{章} & \textbf{书中主题} & \textbf{Semantica 包} & \textbf{当前边界摘要} \\
\midrule
Ch01 & 工程本体论的目的与范围 & \texttt{vol1.ch01} & CQ/案例已入包；概念化边界 runner 与 oracle 缺失 \\
Ch02 & DL、FOL、开放/封闭世界 & \texttt{vol1.ch02} & 受限正向规则与 OWA/CWA 场景可复算；完整 DL、默认/非单调逻辑不支持 \\
Ch03 & CQ、Ontology 101、METHONTOLOGY、OntoClean & \texttt{vol1.ch03} & CQ1 正/单因反例可执行；阶段门禁与 OntoClean 检查器缺失 \\
Ch04 & RDF/RDFS/OWL/SPARQL & \texttt{vol1.ch04} & RDF Dataset 与查询原生；Manchester 解析、章内 shapes、SERVICE 离线执行不完整 \\
Ch05 & DL、SWRL、时态与概率推理 & \texttt{vol1.ch05} & 受限正向规则适配可复算；SWRL built-ins、DL tableau、时态、概率、默认逻辑 fail closed \\
Ch06 & 四类工程应用 & \texttt{vol1.ch06} & 案例文本已迁入；尚无领域 RDF/OWL、SHACL、查询与 exact oracle \\
Ch07 & 图谱流水线、对齐、SHACL & \texttt{vol1.ch07} & SHACL 正/单因反例原生；属性图统一事务与领域 mutation 合同未完成 \\
Ch08 & 幻觉控制、GraphRAG、Text2SPARQL、Agent & \texttt{vol1.ch08} & 规则与案例已入包；动作路由、模型/prompt/快照端到端合同仍阻断 \\
Ch09 & 制造业综合案例 & \texttt{vol1.ch09} & CQ1 精确 oracle 原生；Manchester/SWRL 仅保留，CQ2–CQ7 与端到端推理未假绿 \\
\bottomrule
\end{tabularx}
\end{center}
\color{black}

以上状态以安装版本中的 manifest 为准；书中的表是发布时快照，不替代机器可读清单。

\section*{阅读与复算 / Read and Corroborate}

先读正文，再查看并运行对应包：

\texttt{ONTOLOGY\_ENGINEERING\_ROOT} 指向同时含两卷书源的受控 checkout。另两个环境变量
必须由构建/发布流程写入：一个是所运行 Semantica 的 commit 身份，另一个是该精确
wheel 或运行工件的 SHA-256；空值、占位值或错配值都会 fail closed。

\begin{codebox}
\codeline{semantica\ package\ list\ --volume\ vol1\ --json}
\codeline{semantica\ package\ show\ semantica.chapter\_packages.vol1.ch03\ --json}
\codeline{semantica\ package\ verify-books\ \textbackslash{}}
\codeline{\hspace*{2\ttcw}--book-root\ "\$ONTOLOGY\_ENGINEERING\_ROOT"\ --volume\ vol1\ --json}
\codeline{semantica\ package\ run\ semantica.chapter\_packages.vol1.ch03\ \textbackslash{}}
\codeline{\hspace*{2\ttcw}--runtime-commit\ "\$SEMANTICA\_RUNTIME\_COMMIT"\ \textbackslash{}}
\codeline{\hspace*{2\ttcw}--runtime-artifact-sha256\ "\$SEMANTICA\_RUNTIME\_SHA256"\ \textbackslash{}}
\codeline{\hspace*{2\ttcw}--scenario-id\ OE-V1-CH03-SCN-CQ-ACCEPTANCE-001\ --json}
\codeline{semantica\ package\ verify\ semantica.chapter\_packages.vol1.ch03\ \textbackslash{}}
\codeline{\hspace*{2\ttcw}--runtime-commit\ "\$SEMANTICA\_RUNTIME\_COMMIT"\ \textbackslash{}}
\codeline{\hspace*{2\ttcw}--runtime-artifact-sha256\ "\$SEMANTICA\_RUNTIME\_SHA256"\ \textbackslash{}}
\codeline{\hspace*{2\ttcw}--scenario-id\ OE-V1-CH03-SCN-CQ-ACCEPTANCE-001\ --json}
\end{codebox}
\color{black}

不要在本书目录另装 Jena、owlready2、pySHACL 或 RDFLib 来构造平行答案。
Protégé、Jena、RDF4J、HermiT、Pellet、pySHACL 等仍可作为语言生态与历史实现的
学习对象；但本书所声称的可复算结果只能以 source-locked Semantica 包的运行收据为准。

\section*{权利与责任边界 / Rights and Authority}

本书正文与教学案例仅供学习参考，版权归原书作者所有。外部标准、本体与软件各自
受其许可约束。Semantica 的语义一致性、查询或形状校验结果不等于现实事实、合规
结论、设备操作授权或发布批准；事实权威与决策权威必须由受控来源和具名责任人承担。

\textit{Derived from 《工程本体论》书稿一审0513 正.docx; executable semantics bound to Semantica.}
